\documentclass[tikz,border=10pt]{standalone}
\usepackage[no-math]{fontspec}
\usepackage[brazil]{babel}
\usepackage{amsmath,amssymb}
\IfFileExists{C:/Windows/Fonts/Lato_Regular.ttf}{%
  \setmainfont{Lato}[
    Path = C:/Windows/Fonts/,
    UprightFont = Lato_Regular.ttf,
    BoldFont = Lato_Bold.ttf,
    ItalicFont = Lato_Italic.ttf,
    BoldItalicFont = Lato_BoldItalic.ttf
  ]
  \setsansfont{Lato}[
    Path = C:/Windows/Fonts/,
    UprightFont = Lato_Regular.ttf,
    BoldFont = Lato_Bold.ttf,
    ItalicFont = Lato_Italic.ttf,
    BoldItalicFont = Lato_BoldItalic.ttf
  ]
}{%
  \RequirePackage[default]{lato}
}
\renewcommand{\familydefault}{\sfdefault}

\usepackage{xcolor}
\usepackage{tikz}
\usepackage{pgfplots}
\pgfplotsset{compat=1.18}
\usetikzlibrary{arrows.meta,positioning,shapes.geometric,fit,backgrounds,calc,decorations.pathreplacing,patterns}

% Paleta institucional UEFS / Bioengenharia de Solos
\definecolor{bsNavy}{HTML}{1F4268}
\definecolor{bsNavyDark}{HTML}{152E4A}
\definecolor{bsGold}{HTML}{B38A3D}
\definecolor{bsGoldDark}{HTML}{8C6927}
\definecolor{bsTeal}{HTML}{0E7C7B}
\definecolor{bsForest}{HTML}{2D6A4F}
\definecolor{bsAlert}{HTML}{9E2A2B}
\definecolor{bsEarth}{HTML}{8C5A3C}
\definecolor{bsWater}{HTML}{2E86AB}
\definecolor{bsInk}{HTML}{26313D}
\definecolor{bsMuted}{HTML}{5B6573}

\definecolor{bgNavy}{HTML}{EDF2F7}
\definecolor{bgGold}{HTML}{F8F3E9}
\definecolor{bgTeal}{HTML}{EBF6F5}
\definecolor{bgForest}{HTML}{EEF5F0}
\definecolor{bgAlert}{HTML}{FDF2F2}

\tikzset{
  >=Stealth,
  bsBox/.style={
    draw=#1, fill=white, rounded corners=5pt, line width=1.15pt,
    inner sep=7pt, align=center, font=\footnotesize\sffamily, text=bsInk
  },
  bsTint/.style={
    draw=#1, fill=#1!8!white, rounded corners=5pt, line width=1.15pt,
    inner sep=7pt, align=center, font=\footnotesize\sffamily, text=bsInk
  },
  bsBadge/.style={
    fill=#1, text=white, rounded corners=9pt, inner xsep=8pt, inner ysep=3pt,
    font=\scriptsize\bfseries\sffamily
  },
  bsTitle/.style={font=\normalsize\bfseries\sffamily, text=bsNavy, align=center},
  bsSub/.style={font=\footnotesize\sffamily, text=bsMuted, align=center},
  bsArrow/.style={->, >=Stealth, line width=1.35pt},
  bsDArrow/.style={<->, >=Stealth, line width=1.35pt}
}


\begin{document}
\begin{tikzpicture}[font=\sffamily]

% ============================================================
% TÍTULO GERAL
% ============================================================
\node[bsTitle, anchor=west] at (0, 8.05)
  {Gabião Vivo: Arquitetura Híbrida Pedra–Tela–Ramos Vivos};
\node[bsSub, anchor=west] at (0, 7.60)
  {Composição estrutural da gaiola drenante e inserção de estacas vivas para consolidação radicular};

% ============================================================
% PAINEL ESQUERDO: SEÇÃO DO GABIÃO VIVO
% ============================================================
\begin{scope}[shift={(0,0)}]
  \draw[bsNavy!30, rounded corners=6pt, line width=1.0pt, fill=white]
    (0,-0.55) rectangle (9.65, 7.10);
  \node[bsBadge=bsNavy, rounded corners=5pt, anchor=north west] at (0.30, 6.85) {ARQUITETURA SECCIONAL};

  % Solo retido / reaterro úmido atrás do gabião
  \fill[bsEarth!20]
    (0.45, 0.85) -- (0.45, 5.55) -- (3.15, 5.55) -- (3.15, 1.15) -- (7.75, 1.15) -- (7.75, 0.85) -- cycle;
  % Base / fundação drenante
  \fill[bsNavy!15] (2.95, 0.85) rectangle (6.25, 1.25);
  \draw[bsNavyDark, line width=1.0pt] (2.95, 0.85) rectangle (6.25, 1.25);

  % Nível d'água frontal
  \fill[bsWater!20] (5.55, 1.15) rectangle (9.25, 2.05);
  \draw[bsWater, line width=1.0pt] (5.55, 2.05) -- (9.25, 2.05);
  \node[font=\tiny\bfseries\sffamily, text=bsWater, anchor=south east] at (9.15, 2.08) {N.A. do canal};

  % Geotêxtil filtrante no tardoz do gabião
  \draw[bsNavy, line width=2.4pt, dash pattern=on 5pt off 2.5pt]
    (3.15, 5.45) -- (3.15, 1.25);

  % Gaiolas empilhadas (3 níveis escalonados)
  \foreach \xR/\yB/\yT in {5.65/1.25/2.45, 5.35/2.45/3.65, 5.05/3.65/4.85}{
    \fill[bsEarth!12] (3.25, \yB) rectangle (\xR, \yT);
    \foreach \rx in {3.50, 3.92, 4.34, 4.76}{
      \ifdim \rx pt < \xR pt
        \filldraw[fill=bsNavy!20, draw=bsNavyDark!75, line width=0.6pt]
          (\rx, \yB+0.32) ellipse (0.20 and 0.14);
        \filldraw[fill=bsNavy!26, draw=bsNavyDark!75, line width=0.6pt]
          (\rx+0.12, \yB+0.86) ellipse (0.19 and 0.15);
      \fi
    }
    \draw[bsNavy!55, step=0.30, line width=0.4pt] (3.25, \yB) grid (\xR, \yT);
    \draw[bsNavyDark, line width=1.25pt] (3.25, \yB) rectangle (\xR, \yT);
  }

  % Ramos vivos horizontais inseridos do solo úmido até a face frontal + brotação
  \foreach \y/\xF in {2.45/5.65, 3.65/5.35, 4.85/5.05}{
    \draw[bsEarth!85!black, line width=2.3pt, line cap=round] (1.65, \y+0.18) -- (\xF+0.35, \y-0.02);
    \draw[bsEarth!70!black, line width=0.75pt] (1.65, \y+0.18) -- (1.25, \y+0.34);
    \draw[bsEarth!70!black, line width=0.75pt] (1.65, \y+0.18) -- (1.20, \y-0.02);
    \draw[bsEarth!70!black, line width=0.75pt] (2.25, \y+0.15) -- (1.95, \y-0.10);
    \draw[bsForest, line width=1.6pt, line cap=round] (\xF+0.20, \y) -- (\xF+0.78, \y+0.38);
    \draw[bsForest, line width=1.4pt, line cap=round] (\xF+0.25, \y-0.02) -- (\xF+0.88, \y+0.12);
    \fill[bsForest!85!white] (\xF+0.78, \y+0.38) ellipse (0.20 and 0.10);
    \fill[bsForest!75!white] (\xF+0.88, \y+0.12) ellipse (0.18 and 0.09);
    \fill[bsForest] (\xF+0.56, \y+0.24) ellipse (0.15 and 0.08);
  }

  % Rótulos claros e setas
  \node[font=\scriptsize\bfseries\sffamily, text=bsEarth!90!black, align=left, anchor=west] at (0.55, 5.95)
    {Solo úmido / reaterro\\com enraizamento profundo};

  \node[font=\scriptsize\bfseries\sffamily, text=bsNavy, align=left, anchor=west] (lGeo) at (0.55, 1.65)
    {Geotêxtil\\filtrante};
  \draw[bsArrow, bsNavy, line width=0.95pt] (lGeo.east) -- (3.12, 1.85);

  \node[font=\scriptsize\bfseries\sffamily, text=bsNavyDark, align=left, anchor=west] (lGai) at (5.85, 5.75)
    {Gaiola metálica +\\pedra de mão angular};
  \draw[bsArrow, bsNavyDark, line width=0.95pt] (lGai.west) -- (4.85, 4.55);

  \node[font=\scriptsize\bfseries\sffamily, text=bsForest, align=left, anchor=west] (lRam) at (6.25, 4.15)
    {Ramos vivos horizontais\\brotando na face frontal};
  \draw[bsArrow, bsForest, line width=0.95pt] (lRam.west) -- (5.85, 3.85);

  \node[font=\scriptsize\bfseries\sffamily, text=bsNavyDark, anchor=north] at (4.60, 0.72)
    {Base de fundação drenante};
\end{scope}

% ============================================================
% PAINEL DIREITO: COMPARAÇÃO GABIÃO CONVENCIONAL vs. VIVO
% ============================================================
\begin{scope}[shift={(10.05,0)}]
  \draw[bsGold, rounded corners=6pt, line width=1.2pt, fill=bgGold]
    (0,-0.55) rectangle (5.75, 7.10);
  \node[bsBadge=bsGold, rounded corners=5pt, anchor=north west] at (0.30, 6.85) {ANÁLISE COMPARATIVA};

  \node[bsBox=bsMuted, anchor=north west, text width=4.85cm, align=left] (c1) at (0.30, 6.20)
    {\textbf{\color{bsNavyDark}Gabião Convencional (Mineral)}\\[4pt]
     $\bullet$ Estrutura rígida / estritamente mineral\\[2pt]
     $\bullet$ Dependência contínua da integridade da malha de aço\\[2pt]
     $\bullet$ Baixa integração paisagística e ecológica imediata};

  \node[bsBox=bsForest, anchor=north west, text width=4.85cm, align=left] (c2) at ($(c1.south west)+(0,-0.32)$)
    {\textbf{\color{bsForest}Gabião Vivo (Híbrido)}\\[4pt]
     $\bullet$ Estrutura drenante com matriz solo–pedra–raiz\\[2pt]
     $\bullet$ Enraizamento profundo consolida o talude atrás da gaiola\\[2pt]
     $\bullet$ Integração ecológica superior para recuperação de APP};
\end{scope}

\end{tikzpicture}
\end{document}
